\chapter{工业数据智能案例拆解}

本章不追求给出某个行业的唯一答案，而是展示如何把抽象的模型知识翻译为可执行的问题定义。每个案例都从数据生成、目标、基线、风险和部署反馈五个角度分析。读者应关注决策链条，而不是只记住某个网络名称。

\section{案例一：旋转设备预测性维护}
设设备在时间$t$产生传感器向量$x_t$，维护事件为$y_t$。如果目标是在故障发生前$h$小时发出告警，标签不能简单定义为当前是否故障，而应定义为未来窗口内是否发生故障：
\begin{equation}
 y_t=\mathbb{I}\left(\exists u\in[t+h_1,t+h_2]:\text{failure}(u)=1\right).
\end{equation}
这一定义把预测提前量写进标签。若设备在维护后状态重置，还应将维护事件作为分段边界，避免同一故障周期同时跨越训练和测试集。

\subsection{特征与模型}
短窗口可以提取均值、方差、峰度、峭度、频带能量和包络谱；长窗口可以交给 TCN、LSTM 或 Transformer。树模型常是强基线，因为它能处理非线性阈值和缺失指示变量。序列模型只有在时间顺序包含额外信息时才有必要。

\subsection{评价与告警策略}
样本级 AUPRC 不能直接回答维护人员是否得到足够提前的告警。还应报告事件级召回率、平均提前时间、每台设备每天告警数、重复告警合并后的事件数和误报维护成本。对一个故障周期的多次窗口预测，可以使用首次告警时间定义事件级结果。

\subsection{失效模式}
设备身份泄漏、维护后标签延迟、传感器校准变化和运行工况改变是常见失败原因。模型置信度高不代表告警可信；必须保存触发告警时的输入窗口、模型版本和阈值，以支持复盘。

\section{案例二：视觉表面缺陷检测}
视觉质检的任务粒度可以是图像级分类、目标检测、语义分割或异常检测。选择标准不是“越细越好”，而是缺陷是否需要定位、标注是否足够一致以及生产线是否需要自动剔除。

\subsection{分类、检测与分割}
图像级分类学习$p(y\mid x)$，成本最低但不能指出缺陷位置。检测模型输出边界框和类别，适用于缺陷大小和数量重要的场景。分割模型输出像素级掩码，适用于面积测量或需要精细定位的场景，但标注和评价更昂贵。

交并比为
\begin{equation}
 \operatorname{IoU}(A,B)=\frac{|A\cap B|}{|A\cup B|}.
\end{equation}
检测评价还需指定 IoU 阈值、置信度阈值和类别平均方式。极小缺陷的像素偏差可能显著改变 IoU，因此应补充缺陷召回和漏检成本。

\subsection{异常检测的假设}
当故障样本稀缺时，可以只用正常图像学习重构或特征分布。异常分数可以是重构误差、到正常原型的距离或密度估计结果。该策略隐含“正常数据覆盖足够广”的假设；如果正常工况变化被错误认为异常，系统会产生大量告警。

\subsection{光照与域变化}
相机、镜头、曝光、材料批次和背景变化会产生域偏移。增强策略应模拟真实变化，不能把缺陷边缘随机裁掉后仍宣称任务等价。跨产线测试通常比随机图像划分更能反映部署风险。

\section{案例三：工艺质量回归}
质量指标$y$可能在生产结束后才可测得，模型需要依据中间过程变量预测最终结果。此类任务经常面临标签延迟、批次相关性和控制变量反馈。

\subsection{时间窗口与因果方向}
若预测发生在工艺结束前，只能使用当前时刻之前的变量。把最终检测值或事后调整参数作为特征会产生标签泄漏。对于连续生产过程，应将批次 ID 作为分组变量，同时检查批次内样本是否高度相关。

\subsection{不确定性回归}
点预测不能表示传感器噪声和工艺波动。可以预测均值和方差，使用高斯负对数似然
\begin{equation}
 \mathcal{L}=\frac12\left[\log\sigma_\theta^2(x)+\frac{(y-\mu_\theta(x))^2}{\sigma_\theta^2(x)}\right].
\end{equation}
预测方差较大的样本可进入人工复核，但必须检验方差是否校准：在预测的百分之九十区间内，实际覆盖率应接近百分之九十。

\section{案例四：能源负荷预测}
能源负荷同时具有日周期、周周期、季节变化、节假日效应和突发事件。模型应显式接收日历变量和天气预报，并用滚动时间划分模拟未来预测。

\subsection{基线的重要性}
季节性朴素预测、指数平滑和梯度提升树通常很强。深度模型只有在长历史、多变量或多步预测中显示稳定增益，才值得承担更高部署成本。多步预测可采用直接预测、递归预测或 seq2seq，每种方式都有不同的误差累积特性。

\subsection{峰值风险}
平均 MAE 可能掩盖高峰时段的错误。应按峰值和非峰值分层，并设置峰值低估惩罚。预测结果用于调度时，还需考虑预测分布而非只使用均值。

\section{案例五：工业知识问答与检索增强}
设备手册、维修记录、工艺规范和故障报告构成异构知识库。检索增强生成系统可表示为检索器$R$与生成器$G$的组合：
\begin{equation}
 p(a\mid q,\mathcal{D})=p_G(a\mid q,R(q,\mathcal{D})).
\end{equation}
检索命中不等于回答正确，评价应分离检索召回、证据相关性、答案忠实性和引用完整性。

\subsection{知识库构建}
文档切分应保留章节、设备型号、版本和生效日期。表格、公式和报警代码不能一律当作普通文本处理。每个片段应带来源元数据，回答时引用具体文档位置，便于人工审计。

\subsection{幻觉与拒答}
系统需要允许在证据不足时拒答。可以设置检索分数阈值、证据一致性检查和高风险问题人工审批。对于安全相关操作，不应让语言模型直接执行控制指令。

\section{跨案例的共同验收表}
\begin{table}[htbp]
\centering
\caption{工业数据智能项目验收维度}
\begin{tabular}{p{0.23\textwidth}p{0.63\textwidth}}
\toprule
维度 & 核心问题\\
数据有效性 & 标签如何产生？时间边界是否严格？缺失是否具有含义？\\
模型性能 & 是否超过强基线？是否报告分层和置信区间？\\
鲁棒性 & 换设备、换工况、缺失模态或延迟输入时如何退化？\\
可解释性 & 人员能否追溯输入、证据、版本和决策？\\
系统约束 & 延迟、吞吐、显存、网络和掉线是否满足要求？\\
运营闭环 & 告警如何处理？反馈如何入库？模型何时重新训练？\\
\bottomrule
\end{tabular}
\end{table}

\section{案例练习}
\begin{enumerate}[label=\arabic*.]
 \item 为一个轴承故障任务分别写出样本级和事件级评价指标。
 \item 说明为什么跨产线视觉测试可以发现随机划分无法发现的问题。
 \item 为质量回归模型设计一个预测区间校准实验。
 \item 为工业知识问答系统设计拒答和人工审批规则。
\end{enumerate}

\section{数据治理与标签工程}
工业数据的主要困难往往发生在模型训练之前。一个标签可能来自维修人员填写的工单、实验室检测结果或自动化规则，它们的时间精度、可信度和语义并不相同。标签工程应记录标签来源、生成时间、修订时间、责任人和冲突处理规则。

设同一个事件有多个观测标签$y^{(1)},\ldots,y^{(m)}$，简单多数投票只在标注者质量相近时合理。更稳健的做法是估计标注者可靠性，或把标签冲突作为不确定性的一部分。无论采用何种方法，都不应在看到测试集结果后修改标签定义而不重新冻结评估协议。

\subsection{主数据与实体对齐}
设备名称、传感器编号、产品型号和工单编号常在不同系统中使用不同编码。实体对齐错误会产生重复样本，也可能把一个设备的历史信息泄漏给另一个设备。建议建立稳定的实体主键，并保存原始键到主键的映射版本。删除或合并实体时，要保留可审计的变更记录。

\subsection{质量规则}
数据质量检查至少包含范围检查、单位检查、时间单调性、重复记录、采样间隔、突变率和跨字段一致性。例如压力为负可能是传感器故障，也可能是经过标定的相对压力；规则必须结合测量定义，不能只依据数值直觉。质量规则的失败结果应进入监控，而不是静默丢弃。

\section{在线学习与漂移响应}
当数据分布持续变化时，重新训练不是唯一选择。在线学习可以按到达样本更新参数，窗口学习只使用最近一段数据，周期性再训练则在固定时间点重新构建模型。三种策略的风险不同：在线更新响应快但容易被异常数据污染，窗口学习可能遗忘长期规律，周期性再训练延迟较大但更容易审计。

设新样本到达速率为$\lambda$，单次更新耗时为$c$，若$\lambda c$接近或超过系统可承受上限，就需要采样、批量更新或异步队列。在线系统还需定义更新失败后的回退参数，不能让半完成的更新直接替换生产模型。

\subsection{漂移检测与干预}
漂移检测可针对输入、表示、预测或标签。输入变化可能只是季节性，表示变化可能来自特征工程更新，预测变化可能由阈值改变造成。检测到变化后，应先定位数据管道和业务过程，再决定是否调参、重采样、重新标注或重新训练。自动重训需要设置数据质量门槛和人工审批。

\section{因果边界与决策反馈}
工业模型经常服务于会改变系统状态的决策。模型发出维修告警后，设备被提前维护，未来故障标签因此减少。这种反馈会使“预测性能下降”与“设备真的变健康”难以区分。离线评价应记录干预发生与否，并尽量保留未干预对照。

相关性模型可以预测风险，却不自动回答“采取哪个动作最有效”。若要估计动作效果，需要明确处理变量、结果、混杂因素和时间顺序。不能把观察到的高风险设备被更多维修简单解释为维修导致风险上升。对于高代价动作，应由专家和安全策略共同决定是否允许模型参与。

\section{模型卡片与数据卡片}
模型卡片应说明用途、输入、输出、训练数据范围、已知限制、评价协议、适用设备、失效模式和版本。数据卡片应说明采集时间、覆盖对象、标签规则、缺失模式、许可边界和代表性风险。卡片不是宣传材料，而是让使用者知道哪些结论有证据、哪些结论尚未验证。

\begin{table}[htbp]
\centering
\caption{模型卡片中的风险问题}
\begin{tabular}{p{0.27\textwidth}p{0.59\textwidth}}
\toprule
问题 & 应回答的内容\\
适用范围 & 支持哪些设备、工况、采样率和时间范围？\\
不可用范围 & 哪些输入缺失、漂移或异常时必须拒答？\\
错误代价 & 漏报、误报、延迟和错误自动动作分别造成什么后果？\\
证据边界 & 哪些结论来自测试集，哪些只是经验判断？\\
变更管理 & 数据、模型、阈值和规则如何版本化？\\
\bottomrule
\end{tabular}
\end{table}

\section{案例综合讨论}
同一个工业项目往往同时包含序列、视觉、文本、迁移和持续学习。可以先建立一个能解释、能回退的单模态基线，再逐步加入额外模态和自适应机制。每增加一个模块，都要重新回答输入是否可得、标签是否一致、评价是否公平和失败时如何降级。系统复杂度只有在它解决了已观测到的瓶颈后才值得增加。